\cotefile{161-2}
\scanpage{20}
% archive sheet number: 19

% L1–3
$1^\circ)$ \deleted{\unsure{Théorie vi}} Structure vide sur objets de base $(X_i)_{i\in I}$. \unsure{Représent.ble|Représentables} par un topos $R_0=$ ? Pb 1 ; on a $R_0=\mathrm{Ens.}$ \deleted{\unsure{Ens}} $I$-multi\inserted{\unsure{co|$\omega$}}simpliciaux \inserted{\unsure{opposé-dual}} $={}\in\widehat{\mathbf J}_0$, $\mathbf J_0=$ catégorie \deleted{des} \inserted{familles d'}ensembles finis $(E_i)_{i\in I}$ \unsure{presque tous vide}.
% L2–4, left (slanted)
% slanted
\margin{$J=$ catégorie stricte \unsure{t.q.} mais \unsure{seulemt équivalence} pour $I$}
% L3–4, interline
\[
 I\to\mathbf J_0\to R_0\mathrel{:=}\widehat I\longrightarrow R
\]
$\mathbf J$ est \uline{générateur dans} $R$. \quad $I-\underset{\smile}{\overset{R_0}{J_0}}$

% L4–5
$2^\circ)$ On regarde des flèches \deleted{\unsure{entre des}} entre objets de $R_0$ provenant de $\mathbf J_0$ (celles de $R$), d'où \deleted{$R_{\unsure{1|2|0}}$} \inserted{$J_1$ et $R_1=\rho_\Delta(J_1)$} \deleted{$I\unsure{\to}J_0\unsure{\to}R_0\to R_1\to R$}

% L6–7
$3^\circ)$ on met des axiomes d'égalité sur \deleted{des} \inserted{les} composés (en vertu de $R$) d'où
\[
 \mathbf J\to R_0\to R_1-R_2\to R
\]

% L8
\[
 T_{R_1}(E)=\underline{\mathrm{Hom}}_{\unsure{B|\beta|\Delta}}(R_1,E)\simeq
 \underline{\mathrm{Hom}}(\mathbf J_1,R)\hookleftarrow
 \underline{\mathrm{Hom}}_{\Delta}(R,E)=T_R(E)
\]

% L9–10
où $\mathbf J_1$ est la sous-catégorie pleine de $R$ définie par $b:\mathbf I\to R$, dont $\mathbf R_{2}$ est la \unsure{Ds-enveloppe|$\Delta$-enveloppe} de la petite catégorie $\mathbf J_1$ (cf Pb.1)

% L11–12
\deleted{\unsure{$4^\circ$}} Maintenant \deleted{\unsure{la transformation}} \inserted{le foncteur} $T_R\longrightarrow T_{R_2}$ est pl. fidèle, on sait donc par le N° précédent qu'on a fini par

% L13–20
$4^\circ)$ On met d'autres axiomes \deleted{\illegible{1}} sur la structure. [on peut expliciter lesquels : a) foncteurs \uline{continus} $I_1\unsure{-|\to}E$, i.e. \deleted{satisf} qui transforment certains \deleted{systèmes} cônes inductifs de $I_1$ \inserted{(ceux existant dans $R$)} en cônes exacts \inserted{de $E$} ; \deleted{b) foncteurs} ce sont les foncteurs qui se prolongent en des foncteurs \uline{commutant aux} \uline{lims} $R\to E$ ; plus la condition de transformer certains \deleted{diagrammes de} cônes projectifs finis \uline{de $E$} (ou si on veut, de $R_2'$, le topos cl. \unsure{défini|déf.} par ce système d'axiomes) en cônes exacts ..... . On trouve $I-R_0-R_1-R_2-R_3\simeq R$

% L11–22, left (sideways)
% margin note written sideways
\margin{\diagram{\begin{tikzcd}[ampersand replacement=\&,column sep=small,row sep=small]
 \rho(I) \arrow[d,equal] \& \rho(J_1\unsure{;}J_0) \arrow[d,equal] \& \rho(J_2)\arrow[d,equal] \& {}\\
 R_0\arrow[r] \& R_1\arrow[r,dash] \& R_2\arrow[r,dash]\arrow[d] \& R_3=R\arrow[dl]\\
 J_0\arrow[r]\arrow[u,leftrightarrow] \& J_1\arrow[r]\arrow[u]\arrow[ur,phantom,"\unsure{J_1'}" description]\& J_2 \& {}\\
 I\arrow[u,dash] \& {} \& {} \& {}
 \end{tikzcd}}\deleted{\illegible{1}} $\Delta$-enveloppe, \unsure{pour toutes|pour toutes les} sommes \unsure{provenant de|provenant des} $J_0)=\rho(J'_1)$ $J'_1$ enveloppe \deleted{\unsure{engend|engendré}} par $J_0$ et les flèches données.}

% L21–27
\uline{NB} Comme on a rajouté dans $3^\circ$ et $4^\circ$ des axiomes, on peut les condenser en une seul pas, grâce au fait que $R_1-R_2$ est surjectif sur objets et flèches (cf. Pb 2). Donc on peut \inserted{tout} faire par les trois pas minimum : objets, se \unsure{prendre} flèches \inserted{entre produits finis}, axiomes sur les flèches, exprimables en termes \deleted{de lims} d'inversion de flèches entre \deleted{\unsure{entre}} \uuline{\unsure{paires|termes}} \unsure{formés|formées} entre ces produits \deleted{lims \illegible{3}} \inserted{pour des syst. inductifs de} ou morphismes entre tels formés grâce aux flèches introduites dans les données.

% L25–27, lower left
\margin{\uline{NB} $\widehat{J_1}\to R$ est un \uline{foncteur} de \uline{localisation}}
